% GENERATED FILE. Edit canonical lessons, then run npm run generate:books.
% canonical-source-hash: fnv1a64:41b1e19120615517
% canonical-lessons: GE-C467-scham, GE-C467-schaemen, GE-C467-verlegen, GE-C467-verlegenheit, GE-C467-erroeten

\chapter{Scham und Verlegenheit}
\label{ch:ge-scham-und-verlegenheit}
\hlchaptermodality{467}

\begin{chapteropening}
\textbf{By the end of this chapter:} \emph{I can talk about shame and embarrassment.}

Complete the last lesson of chapter 467: I can talk about shame and embarrassment.
\end{chapteropening}

\section[die Scham]{die Scham — shame}

\label{lesson:GE-C467-scham}



\begin{quote}
Before the new one: say the German for to envy, then the German for to not begrudge.
\end{quote}

\subsection*{You'll want to know: die Scham}
\textbf{die Scham} — \textquotedblleft{}shame\textquotedblright{}.

\begin{quote}
Vor Scham wurde sie ganz rot.
\end{quote}

\begin{quote}
\emph{She went bright red with shame.}
\end{quote}

No plural. \emph{Vor Scham} — with shame.

\begin{cousinweb}
\emph{Scham} is inherited from Old High German \emph{scama}, from Proto-Germanic \emph{skamō} \textquotedblleft{}shame\textquotedblright{} (EtymDB); English \emph{shame} is a cousin.
\end{cousinweb}

\begin{grammarlens}[title={What you've built}]
One more word for telling, explaining and describing at B1.
\end{grammarlens}

\subsection*{Your turn}
\begin{itemize}
\raggedright
  \item \emph{Say it:} \emph{die Scham}
  \item \emph{Say it:} \emph{die Scham}, once more
  \item \emph{Say it:} a sentence of your own with \emph{die Scham}
  \item \emph{Recall:} say the German for to envy, then the German for to not begrudge, then say \emph{die Scham} again
\end{itemize}

\begin{tcolorbox}[breakable,colback=teal!4,colframe=teal!35!black,title={Before you move on}]
What does die Scham mean? (\textquotedblleft{}Shame\textquotedblright{}.) Say it once more.
\end{tcolorbox}
\section[schämen]{schämen — to be ashamed (sich schämen)}

\label{lesson:GE-C467-schaemen}



\begin{quote}
Before the new one: say the German for to not begrudge, then the German for shame.
\end{quote}

\subsection*{You'll want to know: schämen}
\textbf{schämen} — \textquotedblleft{}to be ashamed (sich schämen)\textquotedblright{}.

\begin{quote}
Ich schäme mich für meinen Fehler.
\end{quote}

\begin{quote}
\emph{I'm ashamed of my mistake.}
\end{quote}

Reflexive: \emph{sich schämen für} + accusative, or \emph{sich schämen wegen}.

\begin{cousinweb}
\emph{schämen} comes from Middle High German \emph{schamen} (EtymDB), a verb built on the noun that is today \emph{Scham}, from Proto-Germanic \emph{skamō} \textquotedblleft{}shame\textquotedblright{}.
\end{cousinweb}

\begin{grammarlens}[title={What you've built}]
One more word for telling, explaining and describing at B1.
\end{grammarlens}

\subsection*{Your turn}
\begin{itemize}
\raggedright
  \item \emph{Say it:} \emph{schämen}
  \item \emph{Say it:} \emph{schämen}, once more
  \item \emph{Say it:} a sentence of your own with \emph{schämen}
  \item \emph{Recall:} say the German for to not begrudge, then the German for shame, then say \emph{schämen} again
\end{itemize}

\begin{tcolorbox}[breakable,colback=teal!4,colframe=teal!35!black,title={Before you move on}]
What does schämen mean? (\textquotedblleft{}To be ashamed (sich schämen)\textquotedblright{}.) Say it once more.
\end{tcolorbox}
\section[verlegen]{verlegen — embarrassed, sheepish}

\label{lesson:GE-C467-verlegen}



\begin{quote}
Before the new one: say the German for shame, then the German for to be ashamed.
\end{quote}

\subsection*{You'll want to know: verlegen}
\textbf{verlegen} — \textquotedblleft{}embarrassed, sheepish\textquotedblright{}.

\begin{quote}
Der Junge lächelte verlegen und sagte nichts.
\end{quote}

\begin{quote}
\emph{The boy smiled sheepishly and said nothing.}
\end{quote}

There is also a verb \emph{verlegen} — to misplace, or to move to another date: \emph{Ich habe meine Brille verlegt.}

\begin{cousinweb}
The adjective \emph{verlegen} is usually traced to Middle High German \emph{verlegen}, the past participle of \emph{verligen} \textquotedblleft{}to lie too long, to grow sluggish\textquotedblright{}, built on \emph{liegen} \textquotedblleft{}to lie\textquotedblright{}; EtymDB records only the verb \emph{verlegen} \textquotedblleft{}to misplace\textquotedblright{}.
\end{cousinweb}

\begin{grammarlens}[title={What you've built}]
One more word for telling, explaining and describing at B1.
\end{grammarlens}

\subsection*{Your turn}
\begin{itemize}
\raggedright
  \item \emph{Say it:} \emph{verlegen}
  \item \emph{Say it:} \emph{verlegen}, once more
  \item \emph{Say it:} a sentence of your own with \emph{verlegen}
  \item \emph{Recall:} say the German for shame, then the German for to be ashamed, then say \emph{verlegen} again
\end{itemize}

\begin{tcolorbox}[breakable,colback=teal!4,colframe=teal!35!black,title={Before you move on}]
What does verlegen mean? (\textquotedblleft{}Embarrassed, sheepish\textquotedblright{}.) Say it once more.
\end{tcolorbox}
\section[die Verlegenheit]{die Verlegenheit — embarrassment; awkward situation}

\label{lesson:GE-C467-verlegenheit}



\begin{quote}
Before the new one: say the German for to be ashamed, then the German for embarrassed.
\end{quote}

\subsection*{You'll want to know: die Verlegenheit}
\textbf{die Verlegenheit} — \textquotedblleft{}embarrassment; awkward situation\textquotedblright{}.

\begin{quote}
In seiner Verlegenheit sah er auf den Boden.
\end{quote}

\begin{quote}
\emph{In his embarrassment, he looked at the floor.}
\end{quote}

\emph{Jemanden in Verlegenheit bringen} — to embarrass someone. Plural, for awkward situations: \emph{die Verlegenheiten}.

\begin{cousinweb}
\emph{Verlegenheit} is built on the adjective \emph{verlegen} \textquotedblleft{}embarrassed\textquotedblright{} with \emph{-heit}; \emph{verlegen} is usually traced to a Middle High German participle meaning \textquotedblleft{}having lain too long\textquotedblright{}.
\end{cousinweb}

\begin{grammarlens}[title={What you've built}]
One more word for telling, explaining and describing at B1.
\end{grammarlens}

\subsection*{Your turn}
\begin{itemize}
\raggedright
  \item \emph{Say it:} \emph{die Verlegenheit}
  \item \emph{Say it:} \emph{die Verlegenheit}, once more
  \item \emph{Say it:} a sentence of your own with \emph{die Verlegenheit}
  \item \emph{Recall:} say the German for to be ashamed, then the German for embarrassed, then say \emph{die Verlegenheit} again
\end{itemize}

\begin{tcolorbox}[breakable,colback=teal!4,colframe=teal!35!black,title={Before you move on}]
What does die Verlegenheit mean? (\textquotedblleft{}Embarrassment; awkward situation\textquotedblright{}.) Say it once more.
\end{tcolorbox}
\section[erröten]{erröten — to blush}

\label{lesson:GE-C467-erroeten}



\begin{quote}
Before the new one: say the German for embarrassed, then the German for embarrassment.
\end{quote}

\subsection*{You'll want to know: erröten}
\textbf{erröten} — \textquotedblleft{}to blush\textquotedblright{}.

\begin{quote}
Bei dem Kompliment errötete sie.
\end{quote}

\begin{quote}
\emph{She blushed at the compliment.}
\end{quote}

Perfekt with \emph{sein}: \emph{Sie ist errötet.} A little literary; in speech, \emph{rot werden}.

\begin{cousinweb}
\emph{erröten} is built on \emph{er-} and \emph{röten} \textquotedblleft{}to make red\textquotedblright{}, from \emph{rot} \textquotedblleft{}red\textquotedblright{}, which is inherited from Old High German \emph{rōt}, from Proto-Germanic \emph{raudaz} (EtymDB).
\end{cousinweb}

\begin{grammarlens}[title={What you've built}]
One more word for telling, explaining and describing at B1.
\end{grammarlens}

\subsection*{Your turn}
\begin{itemize}
\raggedright
  \item \emph{Say it:} \emph{erröten}
  \item \emph{Say it:} \emph{erröten}, once more
  \item \emph{Say it:} a sentence of your own with \emph{erröten}
  \item \emph{Recall:} say the German for embarrassed, then the German for embarrassment, then say \emph{erröten} again
\end{itemize}

\begin{tcolorbox}[breakable,colback=teal!4,colframe=teal!35!black,title={Before you move on}]
What does erröten mean? (\textquotedblleft{}To blush\textquotedblright{}.) Say it once more.
\end{tcolorbox}
